% =============================================================================
%  CHAPTER 12: Floating Bodies & Captions
% =============================================================================

\chapter{Floating Bodies \& Captions}

\section{What Are Floats?}

Floats are objects (figures, tables) that \latex{} places at optimal positions in the document, rather than exactly where they appear in the source code.

\begin{table}[H]
\centering
\begin{tabular}{ll}
\toprule
\textbf{Environment} & \textbf{Content} \\
\midrule
\texttt{figure}   & Images, diagrams \\
\texttt{table}    & Tables \\
\texttt{algorithm} & Algorithms (with algorithm package) \\
\texttt{listing}  & Code listings (with listings package) \\
\bottomrule
\end{tabular}
\caption{Float environments}
\end{table}

% ------------------------------------------------------------------
\section{Placement Specifiers}

\begin{table}[H]
\centering
\begin{tabular}{llp{8cm}}
\toprule
\textbf{Specifier} & \textbf{Meaning} & \textbf{Notes} \\
\midrule
\texttt{h} & Here       & Approximately at the current location \\
\texttt{t} & Top        & At the top of the page \\
\texttt{b} & Bottom     & At the bottom of the page \\
\texttt{p} & Page       & On a separate float page \\
\texttt{!} & Override   & Override LaTeX's internal placement rules \\
\texttt{H} & HERE       & Exactly here (requires \texttt{float} package) \\
\bottomrule
\end{tabular}
\caption{Float placement specifiers}
\end{table}

\begin{lstlisting}[language={[LaTeX]TeX}, caption=Float placement]
\begin{figure}[htbp]     % Try here, top, bottom, then page
    ...
\end{figure}

\begin{figure}[H]        % Exactly here (float package)
    ...
\end{figure}

\begin{table}[!ht]       % Override rules, try here then top
    ...
\end{table}
\end{lstlisting}

% ------------------------------------------------------------------
\section{Captions}

\begin{lstlisting}[language={[LaTeX]TeX}, caption=Captions]
% Caption below the content (typical for figures):
\begin{figure}[htbp]
    \centering
    \includegraphics[width=5cm]{image}
    \caption{This is a figure caption.}
    \label{fig:myfigure}
\end{figure}

% Caption above the content (typical for tables):
\begin{table}[htbp]
    \centering
    \caption{This is a table caption.}
    \label{tab:mytable}
    \begin{tabular}{lc}
        ...
    \end{tabular}
\end{table}
\end{lstlisting}

\subsection{Caption Styles (caption package)}

\begin{lstlisting}[language={[LaTeX]TeX}, caption=Caption customization]
\usepackage{caption}

% Change label separator:
\captionsetup{labelsep=period}   % Figure 1. Caption text
\captionsetup{labelsep=quad}     % Figure 1    Caption text
\captionsetup{labelsep=newline}  % Label on separate line
\captionsetup{labelsep=none}     % No label

% Font options:
\captionsetup{font=small}
\captionsetup{font=it}
\captionsetup{labelfont=bf}

% Justification:
\captionsetup{justification=centering}
\captionsetup{justification=raggedright}
\captionsetup{justification=justified}

% Per-environment:
\captionsetup[figure]{font=small, labelfont=bf}
\captionsetup[table]{font=small, skip=5pt}
\end{lstlisting}

% ------------------------------------------------------------------
\section{Side Captions}

\begin{lstlisting}[language={[LaTeX]TeX}, caption=Side captions]
\usepackage{sidecap}

\begin{SCfigure}
    \includegraphics[width=5cm]{image}
    \caption{Caption appears beside the figure.}
\end{SCfigure}
\end{lstlisting}

% ------------------------------------------------------------------
\section{Float Barriers}

Sometimes floats drift too far from their reference. Use barriers to control this:

\begin{lstlisting}[language={[LaTeX]TeX}, caption=Float barriers]
\usepackage{placeins}

% Force all pending floats to be placed before this point:
\FloatBarrier

% Use in a section command:
\makeatletter
\AtBeginSection{\FloatBarrier}
\makeatother
\end{lstlisting}

% ------------------------------------------------------------------
\section{Controlling Float Counters}

\begin{lstlisting}[language={[LaTeX]TeX}, caption=Float counters]
% Maximum floats per page:
\renewcommand{\topfraction}{0.9}     % Max fraction of page for top floats
\renewcommand{\bottomfraction}{0.9}  % Max fraction for bottom floats
\renewcommand{\textfraction}{0.05}   % Min fraction for text
\renewcommand{\floatpagefraction}{0.8} % Min fraction of floats for float page

% Maximum number of floats:
\setcounter{topnumber}{4}            % Max top floats per page
\setcounter{bottomnumber}{4}         % Max bottom floats per page
\setcounter{totalnumber}{10}         % Max total floats per page
\end{lstlisting}
